\begin{proof}[Proof of \cref{obj:thm:whole-null-calibration-resolved}]
Fix an admissible tuple \(v\) and an integer \(n\ge2\). Throughout, write
\[
M:=J_{\mathrm c},
\qquad
B:=\mathfrak b_{n,v}.
\]
These quantities have the values prescribed in \cref{obj:def:explicit-score-ledger}, including its small-sample branch.

\begin{enumerate}
\item First,
\[
B\ge0,
\qquad
M\ge1.
\]
For \(n<4\), these follow from \(B=0\) and \(M=1\). For \(n\ge4\), every summand in the formula for \(B\) is nonnegative: the clipping cutoffs and ranks are positive, and the increment weights \(a_j\) are nonnegative. This also covers an empty increment sum. The coarse rank is a nonnegative integer power of two, hence is at least one.
% lean: CausalSmith.Stat.FinitepHomogeneityDensegamma.ledgerBias_nonneg
% lean: CausalSmith.Stat.FinitepHomogeneityDensegamma.ledgerM_ge_one

\item Let \(P\in\mathcal M_v\). By \cref{obj:lem:original-score-covariance-ledger}, the affine coefficients in
\[
Z_b(c)=Z_{b,0}-cZ_{b,A},
\qquad b\in\{1,2\},
\]
have square-integrable scalar projections, and, for all \(u,z\in\mathbb R^M\),
\[
\begin{aligned}
\operatorname{Var}(u^T Z_{b,0})
&\le \Lambda_{n,v}\|u\|_2^2,\\
\operatorname{Var}(u^T Z_{b,A})
&\le \Lambda_{n,v}\|u\|_2^2,\\
\left|\operatorname{Cov}(u^T Z_{b,0},z^T Z_{b,A})\right|
&\le \Lambda_{n,v}\|u\|_2\|z\|_2.
\end{aligned}
\]
Taking the coordinate directions and summing their second moments proves square integrability of each coefficient vector.

Define the sample event
\[
\mathcal E_{\mathrm{prof}}(P)
:=
\left\{
\text{sample}:
\forall |c|\le\tfrac12,\ 
\left|W(c)-\|\theta_P(c)\|_2^2\right|
\le
128\left(
\sqrt{\Lambda_{n,v}}\|\theta_P(c)\|_2
+\sqrt M\,\Lambda_{n,v}
\right)
\right\}.
\]
By \cref{obj:lem:uniform-affine-profile-bound}, this event is measurable and
\[
P^{\otimes n}\!\left(\mathcal E_{\mathrm{prof}}(P)\right)\ge0.9925.
\]
The same result supplies, when \(\Lambda_{n,v}=0\), the simultaneous almost-sure identity
\[
W(c)=\|\theta_P(c)\|_2^2
\qquad\text{for every }|c|\le\tfrac12.
\]
% lean: CausalSmith.Stat.FinitepHomogeneityDensegamma.original_score_covariance_ledger
% lean: CausalSmith.Stat.FinitepHomogeneityDensegamma.uniform_affine_profile_bound

\item Suppose \(P\in H_0(v)\). By \cref{obj:def:null}, there is a constant \(c_0\) satisfying
\[
c_0\in[-\tfrac12,\tfrac12],
\qquad
\tau_P(x)=c_0\quad\text{for every }x\in[0,1].
\]
If \(n\ge4\), \cref{obj:lem:original-score-mean-ledger} gives
\[
\|\theta_P(c_0)\|_2\le B.
\]
If \(n<4\), the zero-score branch of \cref{obj:def:explicit-score-ledger} gives
\[
\theta_P(c)=0\quad\text{for every }c\in\mathbb R,
\qquad B=0,
\]
so the same bound holds. Since norms are nonnegative and \(c_0\) is an admissible candidate,
\[
\inf_{|c|\le1/2}\|\theta_P(c)\|_2
\le \|\theta_P(c_0)\|_2
\le B.
\]
% lean: CausalSmith.Stat.FinitepHomogeneityDensegamma.whole_null_calibration_resolved

\item We next establish the size bound for \(n\ge4\). The covariance ledger is strictly positive. Indeed, the construction gives
\[
s_n=\lfloor n/2\rfloor\ge2,
\qquad
T_0\ge1,
\qquad
V_0=32T_0^{2-p}>0,
\qquad
L_1\ge16\sqrt{V_0}>0.
\]
Consequently,
\[
\Lambda_{n,v}
=
\frac{L_1^2}{s_n}
+\frac{2L_2^2}{s_n(s_n-1)}
>0.
\]
% lean: CausalSmith.Stat.FinitepHomogeneityDensegamma.ledgerCov_pos

On \(\mathcal E_{\mathrm{prof}}(P)\), define
\[
u_{c_0}:=\|\theta_P(c_0)\|_2.
\]
Then \(0\le u_{c_0}\le B\), and the profiling inequality implies
\[
W(c_0)
\le
u_{c_0}^2
+128\sqrt{\Lambda_{n,v}}\,u_{c_0}
+128\sqrt M\,\Lambda_{n,v}.
\]
Completing the square gives
\[
0\le\left(B-64\sqrt{\Lambda_{n,v}}\right)^2,
\qquad
128\sqrt{\Lambda_{n,v}}\,B
\le B^2+4096\Lambda_{n,v}.
\]
Also, \(M\ge1\) implies
\[
\sqrt M\ge1,
\qquad
\Lambda_{n,v}\le\sqrt M\,\Lambda_{n,v}.
\]
Thus
\[
\begin{aligned}
W(c_0)
&\le B^2+128\sqrt{\Lambda_{n,v}}\,B
       +128\sqrt M\,\Lambda_{n,v}\\
&\le2B^2+4096\Lambda_{n,v}
       +128\sqrt M\,\Lambda_{n,v}\\
&\le2B^2+4224\sqrt M\,\Lambda_{n,v}\\
&\le2B^2+65536\sqrt M\,\Lambda_{n,v}.
\end{aligned}
\]
% lean: CausalSmith.Stat.FinitepHomogeneityDensegamma.profile_null_scalar

\item We justify the comparison with the computed profile minimum and the rounded cutoff. At a fixed sample, define
\[
Z_{b,0}:=Z_b(0),
\qquad
Z_{b,A}:=Z_b(0)-Z_b(1),
\qquad b\in\{1,2\}.
\]
The affine score construction gives \(Z_b(c)=Z_{b,0}-cZ_{b,A}\). Define its quadratic coefficients by
\[
\begin{aligned}
\omega_0&:=\langle Z_{1,0},Z_{2,0}\rangle,\\
\omega_1&:=-\langle Z_{1,A},Z_{2,0}\rangle
           -\langle Z_{1,0},Z_{2,A}\rangle,\\
\omega_2&:=\langle Z_{1,A},Z_{2,A}\rangle.
\end{aligned}
\]
Expansion yields
\[
W(c)=\omega_0+\omega_1c+\omega_2c^2.
\]
Define the smaller endpoint value and the computed profile value by
\[
\omega_{\mathrm{end}}
:=
\min\left\{
\omega_0-\frac{\omega_1}{2}+\frac{\omega_2}{4},
\omega_0+\frac{\omega_1}{2}+\frac{\omega_2}{4}
\right\},
\]
\[
\mathsf W_{\min}
:=
\begin{cases}
\displaystyle
\min\left\{\omega_{\mathrm{end}},
\omega_0-\frac{\omega_1^2}{4\omega_2}\right\},
&
\displaystyle
\omega_2>0,\quad
-\frac12\le-\frac{\omega_1}{2\omega_2}\le\frac12,\\[2mm]
\omega_{\mathrm{end}},&\text{otherwise}.
\end{cases}
\]
The quotient in the first branch is used under \(\omega_2>0\).

For every \(|c|\le1/2\),
\[
\mathsf W_{\min}\le W(c).
\]
To see this in the first branch, complete the square:
\[
W(c)-\left(\omega_0-\frac{\omega_1^2}{4\omega_2}\right)
=
\omega_2\left(c+\frac{\omega_1}{2\omega_2}\right)^2
\ge0.
\]
If \(\omega_2>0\) and the vertex lies to the left of the interval, then \(\omega_1>\omega_2\), and
\[
W(c)-W(-\tfrac12)
=
(c+\tfrac12)\bigl[\omega_1+\omega_2(c-\tfrac12)\bigr]
\ge0.
\]
If the vertex lies to the right, then \(\omega_1<-\omega_2\), and
\[
W(c)-W(\tfrac12)
=
(\tfrac12-c)\bigl[-\omega_1-\omega_2(c+\tfrac12)\bigr]
\ge0.
\]
Finally, if \(\omega_2\le0\), then \(c^2\le1/4\). For \(\omega_1\ge0\),
\[
W(c)-W(-\tfrac12)
=
\omega_1(c+\tfrac12)+\omega_2(c^2-\tfrac14)
\ge0,
\]
whereas for \(\omega_1<0\),
\[
W(c)-W(\tfrac12)
=
\omega_1(c-\tfrac12)+\omega_2(c^2-\tfrac14)
\ge0.
\]
These cases include a constant or linear statistic.
% lean: CausalSmith.Stat.FinitepHomogeneityDensegamma.ledger_profiledMin_le

For the cutoff, define
\[
x_{\mathrm{cut}}
:=2B^2+65536\sqrt M\,\Lambda_{n,v}>0,
\qquad
k_{\mathrm{cut}}
:=\left\lceil\log_2 x_{\mathrm{cut}}\right\rceil.
\]
The defining ceiling inequalities give
\[
\log_2 x_{\mathrm{cut}}
\le k_{\mathrm{cut}}
<\log_2 x_{\mathrm{cut}}+1.
\]
Exponentiating and using the cutoff formula therefore yields
\[
2B^2+65536\sqrt M\,\Lambda_{n,v}
\le h_{n,v}
\le4B^2+131072\sqrt M\,\Lambda_{n,v}.
\]
% lean: CausalSmith.Stat.FinitepHomogeneityDensegamma.ledgerCutoff_bounds

Combining these bounds with the preceding step, on \(\mathcal E_{\mathrm{prof}}(P)\),
\[
\mathsf W_{\min}
\le W(c_0)
\le2B^2+65536\sqrt M\,\Lambda_{n,v}
\le h_{n,v}.
\]
The strict rejection rule in \cref{obj:def:explicit-score-ledger} consequently gives
\[
\psi_{n,v}(\text{sample},\text{seed})=0
\quad\text{on }\mathcal E_{\mathrm{prof}}(P)\times[0,1].
\]

Define the experiment measure and its good event by
\[
\mathbb P_{\mathrm{exp}}:=P^{\otimes n}\otimes\lambda,
\qquad
\mathcal E_{\mathrm{exp}}(P)
:=\mathcal E_{\mathrm{prof}}(P)\times[0,1].
\]
Independence and the unit mass of \(\lambda\) imply
\[
\mathbb P_{\mathrm{exp}}\!\left(\mathcal E_{\mathrm{exp}}(P)\right)
=
P^{\otimes n}\!\left(\mathcal E_{\mathrm{prof}}(P)\right).
\]
The test takes values in \([0,1]\), so pointwise
\[
0\le\psi_{n,v}
\le\mathbf1_{\mathcal E_{\mathrm{exp}}(P)^c}.
\]
Integrating this inequality gives
\[
\begin{aligned}
\mathsf R_n(P,\psi_{n,v})
&\le
\mathbb P_{\mathrm{exp}}\!\left(\mathcal E_{\mathrm{exp}}(P)^c\right)\\
&=1-P^{\otimes n}\!\left(\mathcal E_{\mathrm{prof}}(P)\right)\\
&\le0.0075.
\end{aligned}
\]
For \(n<4\), the zero-test branch gives directly
\[
\mathsf R_n(P,\psi_{n,v})=0\le0.0075.
\]
% lean: CausalSmith.Stat.FinitepHomogeneityDensegamma.integral_le_bad_event
% lean: CausalSmith.Stat.FinitepHomogeneityDensegamma.whole_null_calibration_resolved

\item Now assume
\[
n\ge4,
\qquad
R_{\mathrm{att}}(n,v)<D_v,
\qquad
P\in\mathcal M_v,
\qquad
d(P)\ge R_{\mathrm{att}}(n,v).
\]
Work on \(\mathcal E_{\mathrm{prof}}(P)\). The computed profile value is attained: there exists
\[
c_{\min}\in[-\tfrac12,\tfrac12],
\qquad
W(c_{\min})=\mathsf W_{\min}.
\]
Indeed, the smaller endpoint value is attained at one of the two endpoints. In the vertex branch, choose that endpoint if its value is at most the vertex value; otherwise choose
\[
c_{\mathrm{vert}}:=-\frac{\omega_1}{2\omega_2},
\qquad
W(c_{\mathrm{vert}})
=\omega_0-\frac{\omega_1^2}{4\omega_2}.
\]
Here \(\omega_2>0\) and the branch condition places \(c_{\mathrm{vert}}\) in the candidate interval. In the remaining branch, choose the endpoint attaining \(\omega_{\mathrm{end}}\).
% lean: CausalSmith.Stat.FinitepHomogeneityDensegamma.ledger_profiledMin_attained

Define
\[
u_{\min}:=\|\theta_P(c_{\min})\|_2.
\]
By \cref{obj:lem:original-score-mean-ledger},
\[
u_{\min}
\ge\frac{3}{16}d(P)-16M^{-\gamma}-B.
\]
Substituting the assumed separation and its displayed definition gives
\[
\begin{aligned}
u_{\min}
&\ge
\frac{3}{16}R_{\mathrm{att}}(n,v)-16M^{-\gamma}-B\\
&=4B+1024M^{1/4}\sqrt{\Lambda_{n,v}}.
\end{aligned}
\]

We derive the scalar power bound. Since \(M\ge1\) and \(\Lambda_{n,v}>0\),
\[
M^{1/4}\ge1,
\qquad
\left(M^{1/4}\sqrt{\Lambda_{n,v}}\right)^2
=\sqrt M\,\Lambda_{n,v}>0.
\]
Thus
\[
u_{\min}\ge256\sqrt{\Lambda_{n,v}}\ge0,
\qquad
128\sqrt{\Lambda_{n,v}}\,u_{\min}
\le\frac12u_{\min}^2.
\]
Squaring the signal bound, with all its terms nonnegative, also gives
\[
\begin{aligned}
u_{\min}^2
&\ge
\left(4B+1024M^{1/4}\sqrt{\Lambda_{n,v}}\right)^2\\
&\ge16B^2+1048576\sqrt M\,\Lambda_{n,v}.
\end{aligned}
\]
The lower side of the profiling inequality now yields
\[
\begin{aligned}
W(c_{\min})
&\ge u_{\min}^2
      -128\sqrt{\Lambda_{n,v}}\,u_{\min}
      -128\sqrt M\,\Lambda_{n,v}\\
&\ge\frac12u_{\min}^2-128\sqrt M\,\Lambda_{n,v}\\
&\ge8B^2+524160\sqrt M\,\Lambda_{n,v}\\
&>4B^2+131072\sqrt M\,\Lambda_{n,v}.
\end{aligned}
\]
The last inequality is strict because \(\sqrt M\,\Lambda_{n,v}>0\).
% lean: CausalSmith.Stat.FinitepHomogeneityDensegamma.profile_alternative_scalar

Using the upper cutoff bound and attainment,
\[
h_{n,v}
\le4B^2+131072\sqrt M\,\Lambda_{n,v}
<W(c_{\min})
=\mathsf W_{\min}.
\]
Hence
\[
\psi_{n,v}=1
\quad\text{on }\mathcal E_{\mathrm{exp}}(P),
\qquad
0\le1-\psi_{n,v}
\le\mathbf1_{\mathcal E_{\mathrm{exp}}(P)^c}.
\]
Measurability and boundedness make the test integrable under the probability measure \(\mathbb P_{\mathrm{exp}}\). Integrating the last inequality gives
\[
\begin{aligned}
1-\mathsf R_n(P,\psi_{n,v})
&=\int(1-\psi_{n,v})\,d\mathbb P_{\mathrm{exp}}\\
&\le1-\mathbb P_{\mathrm{exp}}\!\left(\mathcal E_{\mathrm{exp}}(P)\right)\\
&=1-P^{\otimes n}\!\left(\mathcal E_{\mathrm{prof}}(P)\right)\\
&\le0.0075.
\end{aligned}
\]
% lean: CausalSmith.Stat.FinitepHomogeneityDensegamma.integral_le_bad_event
% lean: CausalSmith.Stat.FinitepHomogeneityDensegamma.whole_null_calibration_resolved

\item Finally, if \(n=2\) or \(n=3\), then \(n<4\), and the explicit small-sample branch gives
\[
\psi_{n,v}(\text{sample},\text{seed})=0
\qquad\text{for every sample and seed}.
\]
% lean: CausalSmith.Stat.FinitepHomogeneityDensegamma.ledger_small_sample
\end{enumerate}
\end{proof}
